\subsection{Stokes' Theorem}

\begin{thm}[title=Generalized Stokes' Theorem]
	let $A\subseteq M^{m}$ be a \textit{bounded region} with \textit{boundary} $\p A$.
	then for any $\alpha\in\Omega^{m-1}(M)$ with $\supp(\alpha)\cap A$ cpt, we have:
	\begin{equation*}
		\int_{\p A}\omega \ = \ \int_{A}\d\omega
	\end{equation*}
\end{thm}
remark: $M$ may be a submfld of sth higher-dim!
other remark: we literally havent talked abt bdries lmao

\begin{xmp} \vspace{-0.2in}
	\begin{enumerate}[i)]
		\item sps $\alpha\in\Omega^{2}(\bR^{3})$, $A$ some bounded region in $\bR^{3}$.
			then this yields the divergence theorem.
		\item sps $\alpha\in\Omega^{1}(\bR^{3})$, $A$ some bounded surface in $\bR^{3}$.
			then this yields the classical stokes thm.
		\item sps $\alpha\in\Omega^{0}(M)=C^{\infty}(M)$.
			then this yields the fundamental thm of calculus.
	\end{enumerate}
\end{xmp}
gave 2 exs w pictures. oh well

\begin{crll}
	(more of a special case) if $\alpha$ closed, then:
	\begin{equation*}
		\int_{\p A}\alpha \ = \ 0
	\end{equation*}
\end{crll} \

\begin{defn}
	sps $M$ oriented.
	we say $D\subseteq M$ is a \textbf{region with smooth boundary} if
	$\ex f:C^{\infty}(M)$ with 0 as a reg val and $D=f^{-1}(t\leq0)$.
	we write $\p D=f^{-1}(0)$ and $\trm{int}(D)=f^{-1}(t<0)$.
\end{defn}
fix a special atlas covering $D$:
\begin{itemize}
	\item cvr $\p D$ by submfld charts $(U,\vphi)$ compat w orientation of $M$, ie:
		\begin{equation*}
			\vphi(U\cap\p D) \ = \ \vphi(U)\cap\set{x_{1}=0} \qquad
			\vphi(U\cap\trm{int}D) \ = \ \vphi(U)\cap\set{x_{1}<0}
		\end{equation*}
		after possibly flipping sign of $x_{1}$
	\item cvr interior by more oriented charts
\end{itemize}
this gives an oriented atlas on $D$.

\begin{prop}
	when restricting to $\set{x_{1}=0}$, we get an oriented atlas on $\p D$
\end{prop} \

\begin{pf}
	let $\tau:\vphi_{i}(U_{i}\cap U_{j})\gto\vphi_{j}(U_{i}\cap U_{j})$ be
	the transition map.
	given $x=(0,x_{2},\dots,x_{m})$, consider $T_{x}\tau$:
	the top left is $>0$ since ``away from $D$" is preserved, the rest of
	the first row is 0 [cz bdry is preserved? idk], and ignoring the
	first column, the remaining matrix is made up of the jacobian of
	the transition map btwn charts on $\p D$.

	we know $\det T_{x}\tau>0$, so $\det$(that other thing) is +ve, so the atlas
	$\left(U\big\rvert_{\p D},\vphi\big\rvert_{\p D}\right)$ is oriented.
\end{pf}
note: the orientation we put on $\p D$ is engineered to avoid a minus sign.

\lecdate{Lec 22 - Mar 31 (Week 12)}
today, we prove stokes.

recall we can pick oriented atlas on $D$ which induces an orientation on $\p D$.
since $\alpha$ cpt supp on $D$, its sppted on finitely many charts.
pick a partition of unity $\rho_{i}$ sppted on:
\begin{equation*}
	\set{U_{i}} \cup (M\sm D\cap\supp(\alpha))
\end{equation*}
then, we can write:
\begin{equation*}
	\alpha \ = \ \sum_{i}\rho_{i}\alpha \qquad
	\rho_{i}\alpha \trm{ supp on } U_{i}
\end{equation*}
thus, its enough to show stokes for $\alpha$ cpt supp on $\bR^{m}$, $D=\set{x:x_{1}\leq0}$.
then:
\begin{equation*}
	\alpha \ = \ \sum_{i=1}^{m}f_{i}\di x_{1}\wedge\cdots\wedge\hat{\d x_{i}}\wedge\cdots\wedge\d x_{m}
\end{equation*}
evaluating at $(\p_{2},\dots,\p_{m})$, we get:
\begin{equation*}
	\int_{\set{x_{1}=0}}\alpha \ = \ \int_{\bR^{m-1}}\iota^{*}\alpha
	\ = \ \int_{\bR^{m-1}}f_{1}\di x_{2}\wedge\cdots\wedge\d x_{m}
\end{equation*}
otoh:
\begin{equation*}
	\d\alpha \ = \ \left[\sum_{i=1}^{m}(-1)^{i+1}\p_{i}f_{i}\right]
	\d x_{1}\wedge\cdots\wedge\d x_{m}
\end{equation*}
so:
\begin{equation*}
	\int_{\set{x_{1}\leq0}}\d\alpha
	\ = \ \sum_{i=1}^{m}\int_{-\infty}^{\infty}\cdots \int_{-\infty}^{\infty}\int_{-\infty}^{0}
	\p_{i}f_{i}\di x_{1}\wedge\cdots\wedge\d x_{m}
\end{equation*}
note that when $i\neq1$, this is equal to 0.
indeed, we can integrate $\d x_{i}$ first via fubinis, then by ftc:
\begin{equation*}
	\int_{-\infty}^{\infty}\p_{i}f_{i} \ = \
	\lim_{x\sto\infty}f_{i}-\lim_{x\sto-\infty}f_{i} \ = \ 0
\end{equation*}
when $i=1$, we get:
\begin{equation*}
	\int_{\bR^{m-1}}\left[\int_{-\infty}^{0}\p_{1}f_{1}\di x_{1}\right]
	\ = \ \int_{\bR^{m-1}}f_{1}(0,x_{2},\dots,x_{m})\di x_{2}\cdots\d x_{m} \ = \ 0
\end{equation*}

\begin{crll}
	if $M^{m}$ orientable, $\alpha\in\Omega^{m}(M)$ (w cpt sppt) is exact, ie $\alpha=\d\omega$,
	and $\omega$ cpt sppt, then $\int_{M}\alpha=0$
\end{crll}
fact: if $M$ is connected, the converse is also true. 

\begin{crll}
	if $M$ orientable cpt conn, then $H^{m}(M)\simeq\bR$, where $H^{m}$ is the
	$m$-th de Rham cohomology group.

	if $M$ not cpt, then $H_{c}^{m}(M)\simeq\bR$, where we restrict to forms
	w cpt sppt.
\end{crll} \

\begin{prop}
	every orient $M$ admits a \textit{volume form}, ie an $m$-form vanishing nowhere
\end{prop}
note: if $M$ cpt, $\int_{M}\omega\neq0$.
otw, may not have finite integral.

\begin{pf}
	take locally finite oriented atlas, partition of unity $\rho_{i}$, and:
	\begin{equation*}
		\omega \ = \ \sum_{i}\vphi_{i}^{*}(\rho_{i}\circ\vphi_{i}^{-1})
		\d x_{1}\wedge\cdots\wedge\d x_{m}
	\end{equation*}
	note each summand extends over $M$ since its 0 outside cpt subset of $U_{i}$.
	an orientation ensures signs are all the same, so $\omega\neq0$.
\end{pf}

\newpage
\begin{prop}
	if $M$ admits a vol form, its orientable
\end{prop} \

\begin{pf}
	at each pt, given basis $X_{1},\dots,X_{m}$ of $T_{p}M$,
	$\omega_{p}(X_{1},\dots,X_{m})>0$ or $<0$.
	this defns an orientation at a pt; check that this defns a
	global orientation [see hw].
\end{pf} \

\begin{thm}
	sps $\omega\in\Omega^{k}(M)$ closed, $F:\bR\times S\gto M$ smooth,
	considered as a family of smooth $F_{t}:S\gto M$.
	then $\int_{S}F_{t}^{*}\omega$ is indep of $t$.
\end{thm} \

\begin{pf}
	consider $D=[a,b]\times S$:
	\begin{equation*}
		0 \ = \ \int_{D}F^{*}\d\omega \ = \ \int_{D}\d F^{*}\omega
		\ = \ \int_{\p D}F^{*}\omega \ = \ \int_{S}F_{b}^{*}\omega-\int_{S}F_{a}^{*}\omega
	\end{equation*}
\end{pf}
